%
% Erstellt von Daniel Falkner
% daniel.falkner@akad.de
% 
\documentclass[xcolor=dvipsnames]{beamer}
\usepackage[T1]{fontenc}
\usepackage[utf8]{inputenc}
\usepackage[justification=centering,figurename=Abb.]{caption}

\usetheme{Warsaw}
\usecolortheme[named=OliveGreen]{structure}
\renewcommand\thempfootnote{\arabic{mpfootnote}}

\newcommand*{\Title}{Warum ist ein neues Internetprotokoll notwendig?} %Titel
\subtitle{Modul WIN11} %Untertitel
\newcommand*{\Author}{Daniel Falkner + Eugen Grinschuk} %Name
\institute{AKAD Pinneberg + Stuttgart} %Uni
\titlegraphic{\includegraphics[scale=0.2]{akad_logo.png}} %Logo

\title{\Title}
\author{\Author}
\date{\today}

%Pdf Metainformationen
\subject{\Title}
\keywords{}

\begin{document}

%Titelseite
\begin{frame}
    \titlepage
\end{frame}

%Logo auf allen weiteren Folien
%\logo{\includegraphics[scale=0.1]{akad_logo.png}}

%Inhaltsverzeichniss
\frame{\tableofcontents} 

%Ab hier die Folien 
 
\section{Hintergründe}
\begin{frame} %%Eine Folie
 \frametitle{Allgemein}
Das Internet Protocol (IP) ist die Grundlage für die Datenübertragung im Internet. Der sendende Rechner entlässt ein Paket ins Netz. Router, welche die verschiedenen Teilnetze verbinden, sorgen dafür, dass es auf einem (nicht von vornherein festgelegten) Weg durch die Teilnetze den Empfänger erreicht. Sender und Empfänger sind durch eindeutige IP-Adressen identifiziert, die den eigentlichen Nutzdaten vorangestellt werden.
\end{frame}


\subsection{IPv4}
\begin{frame}
 \frametitle{Internet Protocol Version 4}
  \begin{block}{IPv4}
	  \begin{itemize}
	  	\item Definiert 1981 im RFC \footnote{Request for Comments} 791
	  	\item Beispiel Adresse: 8.8.8.8
  		\item nur $2^{32}$ (4294967296) = 4 Milliarden Adressen
	  \end{itemize}
  \end{block}
\end{frame}

\subsection{IPv6}
\begin{frame}
  \frametitle{Internet Protocol Version 6}
  \begin{block}{IPv6}
	  \begin{itemize}
  		\item Definiert 1998 im RFC 2460 
  		\item Beispiel Adresse: 2001:4860:4860::8888
	  \end{itemize}
  \end{block}
\end{frame}


%\frame{\tableofcontents[currentsection]}

\section{Nachteile IPv6}
\begin{frame}
  \begin{alertblock}{Nachteile Internet Protocol Version 6}
	  \begin{itemize}
	    \item Software muss IPv6 fähig sein. (Kosten)
	    \item Mögliche Sicherheitsprobleme durch Wegfallen von IPv4 NAT \footnote{Network Address Translation}
	  \end{itemize}
  \end{alertblock}
\end{frame}

\section{Vorteile IPv6}
\begin{frame}
  \begin{block}{Vorteile Internet Protocol Version 6}
	  \begin{itemize}
  		\item $2^{128}$ (3,402823669 * $10^{38}$) = > 340 Sextillionen Adressen
  		\item Jedes Endgerät hat eine eindeutige IP-Adresse (kein NAT \footnote{Network Address Translation})
  		\item Autokonfiguration SLAAC \footnote{Stateless Address Autoconfiguration} mit DAD \footnote{Duplicate Address Detection}
 
	  \end{itemize}
  \end{block}
\end{frame}



\section{Fazit und Ausblick}
\begin{frame}
 \frametitle{Lorem ipsum}
  Lorem ipsum dolor sit amet, consetetur sadipscing elitr, sed diam nonumy eirmod tempor invidunt ut labore et dolore magna aliquyam erat, sed diam voluptua. At vero eos et accusam et justo duo dolores et ea rebum. Stet clita kasd gubergren, no sea takimata sanctus est Lorem ipsum dolor sit amet.
\end{frame}

\section{Quellen}
\begin{frame}
 \frametitle{Quellen}
  \begin{block}{}
	  \begin{itemize}
  		\item IPv4 \url{http://tools.ietf.org/html/rfc791}
  		\item IPv6 \url{http://tools.ietf.org/html/rfc2460}
  		\item IPv6 SLAAC \footnote{Stateless Address Autoconfiguration} \url{http://tools.ietf.org/html/rfc4862}
	  \end{itemize}
  \end{block}
\end{frame}

\subsection*{Ende}
\begin{frame}
	\begin{block}{}	
		\begin{center}
			Vielen Dank für Ihre Aufmerksamkeit. \\
			\Author{}
		\end{center}	
	\end{block}
\end{frame}

\end{document}


